\documentclass[12pt]{article}
\usepackage[T1]{fontenc}
\usepackage[utf8]{inputenc}
\usepackage{lmodern}
\usepackage{textcomp}
\usepackage{enumitem}
\usepackage{geometry}
\geometry{margin=1in}
\begin{document}
\begin{center}
Answer Key - Geography - K-12 Level Assessment 17 \
\small (Answers are ordered exactly as in the question paper.)
\end{center}

\section*{Multiple Choice}
\begin{enumerate}[label=\arabic*.]
\item \textbf{Answer:} C
\\ \textit{Options:} A) push factors.; B) pull factors.; C) distance decay.; D) migration selectivity.
\\ \textit{Note:} Distance decay refers to the principle that the likelihood of migration decreases as the distance between the origin and destination increases, reflecting how geographic separation affects social and economic interactions.
\item \textbf{Answer:} B
\\ \textit{Options:} A) First Agricultural Revolution.; B) Second Agricultural Revolution.; C) Third Agricultural Revolution.; D) Fourth Agricultural Revolution.
\\ \textit{Note:} The emergence of an urban industrial workforce in Europe during the Second Agricultural Revolution facilitated increased agricultural productivity, which in turn supported urbanization and the growth of industry by providing both a labor force and a surplus of food.
\item \textbf{Answer:} C
\\ \textit{Options:} A) A central park; B) Prominent religious buildings; C) A large central square surrounded by government and administrative buildings; D) Luxury apartment buildings
\\ \textit{Note:} A large central square surrounded by government and administrative buildings is a distinct characteristic of East European cities, reflecting their historical development as centers of governance and public life.
\item \textbf{Answer:} B
\\ \textit{Options:} A) Europe and North Africa; B) Asia and Latin America; C) Sub-Saharan Africa and Oceania; D) Central Europe and Australia
\\ \textit{Note:} The correct answer is Asia and Latin America because these regions have been the primary sources of immigrants to the United States since 1980, driven by factors such as economic opportunities, political instability, and family reunification.
\item \textbf{Answer:} D
\\ \textit{Options:} A) Congestion in the core; B) High cost of labor in the core; C) High density in the core; D) More infrastructure in the core
\\ \textit{Note:} The correct answer is "More infrastructure in the core" because, in the core-periphery model, spread effects refer to positive impacts and developments that radiate outward from the core to the periphery, while increased infrastructure in the core itself does not represent a spread effect but rather a concentration of resources.
\end{enumerate}

\section*{True / False}
\begin{enumerate}[label=\arabic*.]
\setcounter{enumi}{5}
\item \textbf{Answer:} False
\\ \textit{Options:} A) True; B) False
\\ \textit{Note:} A software engineer is part of the basic employment sector because their work in technology development contributes to the foundational infrastructure that supports various industries and economic activities.
\item \textbf{Answer:} True
\\ \textit{Options:} A) True; B) False
\\ \textit{Note:} The statement is true because the Romance language family originated from Latin, the language of the Roman Empire, and includes Italian, Spanish, Portuguese, and Romanian as they all evolved from this common linguistic root.
\item \textbf{Answer:} False
\\ \textit{Options:} A) True; B) False
\\ \textit{Note:} Freeways often promote urban sprawl by making it easier for people to live farther away from city centers, leading to the development of suburban areas rather than encouraging compact city growth.
\item \textbf{Answer:} True
\\ \textit{Options:} A) True; B) False
\\ \textit{Note:} Japan has implemented various measures, such as promoting the use of traditional Japanese in education and media, to limit the influence of English and other languages on its native language and preserve its cultural identity.
\item \textbf{Answer:} False
\\ \textit{Options:} A) True; B) False
\\ \textit{Note:} The area surrounding an urban center that is economically dependent is called its hinterland, not its threshold, which refers to the minimum market size needed to support a business.
\end{enumerate}

\section*{Long Form Response}
\begin{enumerate}[label=\arabic*.]
\setcounter{enumi}{10}
\item \textbf{Answer:} The Shiite branch of Islam, particularly in Iran, holds a unique belief in the authority of Imams who are considered to be the rightful leaders descended from Ali, the cousin and son-in-law of the Prophet Muhammad. Shiites believe that these Imams are divinely appointed and possess special spiritual and temporal authority, which sets them apart from other Islamic traditions. This belief significantly influences Iran's political and religious landscape, as the country is predominantly Shiite and its governance is intertwined with religious leadership. The historical significance of Shiite Islam in Iran is evident in the Islamic Revolution of 1979, which emphasized the role of the Imams and established a theocratic system based on Shiite principles. Overall, Shiite beliefs shape the cultural and political identity of Iran, making it a focal point for Shiite Muslims worldwide.
\\ \textit{Note:} correct because it accurately describes the Shiite belief in the divinely appointed authority of Imams descended from Ali, emphasizing their influence on Iran's predominantly Shiite political and religious landscape.
\item \textbf{Answer:} Proximity to markets is crucial for many industries because it reduces transportation costs and ensures timely delivery of products. Industries that produce perishable goods, such as food, benefit greatly from being close to their markets to maintain freshness and meet consumer demand. However, bulk-reducing products, like metals or paper, do not require closeness to their markets because they are processed from raw materials that are much heavier than the final product. Since these products lose bulk during production, it is often more economical to locate processing plants near the raw material sources rather than the final consumers. This allows companies to minimize shipping costs for the heavier raw materials while still delivering the lighter finished products to markets effectively.
\\ \textit{Note:} correct because it highlights how proximity to markets minimizes transportation costs and enhances delivery efficiency for perishable goods, while explaining that bulk-reducing products, which are lighter than their raw materials, can be transported over longer distances without significant cost penalties.
\end{enumerate}

\vfill
\noindent\textit{End of Answer Key}
\end{document}